
\chapter{Sistemas Embarcados}
\label{chap:embarcados}

% Apresentar o sistema embarcado do DoseCerta e sua função
% no controle do dispensador.

\section{Plataforma de Desenvolvimento}
\label{sec:plataforma-embarcada}

% Identificar a placa, o microcontrolador ou computador embarcado.
% Informar características relevantes para o projeto e justificar
% a escolha com base nos requisitos.

\section{Ambiente de Desenvolvimento}
\label{sec:ambiente-embarcado}

% Informar linguagem de programação, bibliotecas, ferramentas
% e ambiente utilizados, conforme o desenvolvimento realizado.

\section{Organização do Firmware}
\label{sec:organizacao-firmware}

% Explicar como o código está organizado:
% módulos, funções principais e responsabilidades.
% Inserir diagrama ou estrutura de arquivos, se pertinente.

\section{Inicialização e Configuração}
\label{sec:inicializacao}

% Descrever a sequência de inicialização do sistema,
% incluindo configuração de periféricos e verificação inicial,
% quando implementadas.

\section{Controle do Mecanismo de Dispensação}
\label{sec:controle-dispensacao}

% Explicar a lógica de controle do mecanismo:
% comandos, sequência de operação, condições de parada
% e interação com sensores, se aplicável.
% Relacionar ao projeto mecânico e eletrônico.

\section{Leitura de Sensores e Acionamento de Atuadores}
\label{sec:leitura-acionamento}

% Descrever como o firmware recebe sinais dos sensores
% e controla os atuadores.
% Informar tratamento de sinais, temporização ou calibração
% apenas quando fizerem parte da implementação.

\section{Temporização e Lógica de Alertas}
\label{sec:temporizacao-alertas}

% Descrever como os horários de dispensação são tratados
% e como os alertas são acionados, caso essas funções
% estejam implementadas no sistema embarcado.

\section{Comunicação com Outros Módulos}
\label{sec:comunicacao-embarcada}

% Explicar as interfaces de comunicação efetivamente utilizadas
% para integração com backend, aplicativo ou outros módulos.
% Especificar mensagens, comandos ou protocolos, se pertinente.

\section{Tratamento de Falhas}
\label{sec:tratamento-falhas-embarcado}

% Descrever as falhas previstas e o comportamento do firmware
% diante delas, como leituras inválidas, ausência de resposta
% ou interrupção da operação.
% Não declarar mecanismos de segurança que não foram implementados.

\section{Fluxo de Execução}
\label{sec:fluxo-execucao}

% Apresentar o fluxo principal do firmware.
% Inserir fluxograma ou pseudocódigo, se disponível.

\section{Síntese da Implementação Embarcada}
\label{sec:sintese-embarcada}

% Resumir as funções implementadas e as limitações atuais
% do sistema embarcado.
